\bookchapter{Bab 18: Sandaran Data \& Pemindahan}

\chapterintro{: Foto video (jika penyegerakan awan tidak diaktifkan), nota, SMS, rekod panggilan, fail yang dimuat turun.}

\booksection{Disimpan di pelayan App, log masuk semula pada telefon pintar baru dan ada semula}

: Senarai rakan WeChat, pesanan Taobao, rekod transaksi dalam App perbankan. Perkara-perkara ini disimpan di pelayan syarikat berkenaan, bukan di telefon pintar anda.

\begin{itemize}
  \item \textbf{Dalam kawasan kelabu, memerlukan operasi khusus}: Rekod sembang WeChat (walaupun pelayan WeChat menyimpannya, anda sendiri tidak boleh memuat turun dari pelayan), kredensial peranti yang terikat dengan perbankan telefon pintar, kunci dalam App pengesahan dua langkah (seperti Google Authenticator).
  \item \textbf{Jadi yang perlu anda lindungi dengan utama ialah kategori pertama dan ketiga.}2. Sandaran Seluruh Peranti: Selesai Semua Sekali
  \item \textbf{Sandaran seluruh peranti ibarat mengambil "gambar segera" lengkap telefon pintar — tetapan sistem semasa, susun atur desktop, data App, foto, kenalan, semuanya dibungkus.}Pelan 1: Sandaran tempatan komputer (paling boleh dipercayai, disyorkan untuk mereka yang ada komputer)
\end{itemize}

iPhone

\booksection{: Sambungkan telefon pintar ke komputer dengan kabel data, buka iTunes (Windows versi lama) atau Finder (Mac), dalam halaman peranti pilih “Sandar Sekarang”. Perhatikan bahawa anda perlu menanda “Sulitkan sandaran”, jika tidak kata laluan dan data kesihatan tidak akan disandarkan.}

Android

\booksubsection{: Setiap jenama mempunyai perisian komputer sendiri, seperti Pembantu Telefon Huawei, Pembantu Xiaomi, Samsung Smart Switch. Sambungkan ke komputer dan ikut arahan.}

\begin{itemize}
  \item \textbf{Kelebihan terbesar sandaran awan ialah "tidak perlu anda risau", selesai secara automatik. Kelemahan terbesar ialah pemulihan memerlukan masa muat turun yang lama, dan ruang percuma biasanya tidak mencukupi.}Pelan 2: Sandaran awan (paling menjimatkan masa, disyorkan untuk semua orang)
  \item \textbf{iOS}: Tetapan Nama anda iCloud Sandaran Awan iCloud Hidupkan. Apabila telefon pintar dipasang ke sumber kuasa, disambungkan ke Wi-Fi, dan skrin dikunci, ia akan sandar secara automatik. 5GB percuma kebanyakannya tidak mencukupi, disyorkan belanjakan 6 yuan untuk menaik taraf ke 50GB.
  \item Android
\end{itemize}

\booksubsection{: Awan setiap jenama mempunyai fungsi sandaran. Telefon pintar Huawei pergi ke tetapan cari “sandaran awan”, Xiaomi pergi ke “Perkhidmatan Awan Xiaomi”, OPPO/vivo juga serupa. Selain itu ada sandaran terbina dalam Google (Tetapan Sistem Sandaran), hidupkan juga.}

\begin{itemize}
  \item \textbf{Kelebihan terbesar sandaran awan ialah "tidak perlu anda risau", selesai secara automatik. Kelemahan terbesar ialah pemulihan memerlukan masa muat turun yang lama, dan ruang percuma biasanya tidak mencukupi.}3. Strategi Sandaran Berkategori: Perkara Penting Simpan Satu Salinan Tambahan Secara Berasingan
  \item \textbf{iOS}Foto dan Video (Aset Digital Paling Berharga)
  \item Pelan insurans berganda
\end{itemize}

\booksection{: Aktifkan penyegerakan automatik foto iCloud/Google Foto, pada masa yang sama setiap beberapa bulan salin foto secara manual ke komputer atau cakera keras luaran. Dengan ini satu salinan di awan, satu salinan tempatan, saling menjadi sandaran.}

Pelan tanpa belanja

\booksubsection{: Google Foto memberikan 15GB secara percuma, boleh ditetapkan kepada “kualiti dimampatkan” untuk menjimatkan ruang. Atau setiap sebulan gunakan kabel data untuk menyalin semua foto ke komputer.}

\begin{itemize}
  \item \textbf{Kenalan (Mudah Diabaikan tetapi Sangat Penting)}Kenalan iPhone disegerakkan dengan iCloud secara lalai, selagi penyegerakan diaktifkan, sandaran tambahan tidak diperlukan. Tetapi jika anda mahu lebih selamat, boleh pergi ke iCloud.com mengeksport fail vCard dan simpan di komputer.
  \item \textbf{Kenalan Android biasanya disegerakkan dengan Google atau awan jenama. Pergi ke tetapan App Kenalan untuk mengeksport fail .vcf, itu adalah satu sandaran lengkap.}Rekod Sembang WeChat (Urus secara berasingan, yang paling penting)
\end{itemize}

\booksubsection{Ini mesti dibincangkan secara berasingan, kerana rekod sembang WeChat dalam sandaran seluruh peranti sering kali tidak lengkap, dan alat pertukaran telefon pintar sama sekali tidak boleh memindahkannya.}

\begin{itemize}
  \item Kaedah 1: Pindah ke telefon pintar baru
  \item . WeChat Saya Tetapan Umum Pemindahan \& Sandaran Rekod Sembang Pindah Pindah ke telefon pintar/tablet. Telefon pintar lama dan baru sambung ke Wi-Fi yang sama, imbas kod dan boleh dipindahkan. Perhatian: Sebelum pemindahan selesai, jangan log keluar WeChat dari telefon pintar lama.
\end{itemize}

\booksubsection{Kaedah 2: Sandar ke komputer}

. Pasang WeChat versi PC di komputer, tiga garis di sudut kiri bawah Pemindahan \& Sandaran Sandaran \& Pemulihan Sandar rekod sembang ke komputer. Selepas itu apabila bertukar telefon pintar, boleh dipulihkan ke telefon pintar baru.

\begin{itemize}
  \item \textbf{Kaedah 3: Sandar ke awan}. WeChat juga menyokong sandaran ke iCloud atau ruang awan telefon pintar, tetapi perhatian: sandaran ini hanya boleh dipulihkan pada telefon pintar semasa anda, tidak boleh dipulihkan merentas peranti. Jika bertukar telefon pintar, masih perlu menggunakan dua kaedah pertama.
  \item \textbf{4. Panduan Pertukaran Telefon Pintar: Dari Telefon Pintar Lama ke Telefon Pintar Baru}Bertukar telefon pintar baru adalah perkara gembira, tetapi jika data tidak dipindahkan, ia menjadi sakit kepala. Tiga situasi diurus secara berasingan:
  \item \textbf{Pertukaran jenama sama (paling mudah)}iPhone iPhone baru
\end{itemize}

\booksection{: Semasa mengaktifkan, letakkan telefon pintar lama di sebelahnya, “Mula Cepat” akan muncul secara automatik, imbas corak pada telefon pintar baru dengan telefon pintar lama. Semua data, tetapan, susun atur App akan dipindahkan tanpa perubahan, termasuk kata laluan Wi-Fi juga akan dibawa secara automatik. Kira-kira memerlukan 20-30 minit.}

Android jenama sama

\booksubsection{: Gunakan alat pertukaran telefon pintar terbina dalam jenama. Huawei gunakan “Klon Telefon”, Xiaomi gunakan “Pertukaran Telefon Xiaomi”, OPPO gunakan “Perpindahan Telefon”, vivo gunakan “Hantaran Bersama”. Langkahnya hampir sama: Pasang App yang sama pada telefon pintar lama dan baru, imbas kod untuk menyambung, pilih item yang ingin dipindahkan, mula pemindahan.}

\begin{itemize}
  \item \textbf{Pertukaran merentas jenama (Android ke Android, jenama berbeza)}Gunakan alat pertukaran telefon pintar pihak ketiga atau gunakan App pertukaran terbina dalam jenama lama (kebanyakan App pertukaran jenama juga menyokong pertukaran ke telefon pintar Android lain). Yang boleh dipindahkan: Kenalan, foto, SMS, rekod panggilan, sebahagian App. Yang tidak boleh dipindahkan: Status log masuk dalam App (anda perlu log masuk semula), rekod sembang WeChat (gunakan fungsi pemindahan WeChat sendiri).
  \item \textbf{Pertukaran merentas platform (Android ke iPhone, atau sebaliknya)}Android iPhone
\end{itemize}

\booksubsection{: Apple mempunyai App “Pindah ke iOS”, muat turun di telefon pintar Android, semasa mengaktifkan iPhone baru pilih “Pindahkan Data dari Peranti Android”. Boleh memindahkan kenalan, foto, SMS dan data asas lain, tetapi data App tidak boleh dipindahkan. Rekod sembang WeChat masih perlu menggunakan fungsi pemindahan WeChat sendiri.}

iPhone Android

\booksubsection{: Setiap jenama Android mempunyai alat pemindahan yang sepadan, tetapi jujur, pengalaman pemindahan merentas platform tidak sebaik platform yang sama. Perkara paling penting buat terlebih dahulu: Eksport kenalan, sandar foto, pindah rekod sembang WeChat. Baki boleh dibuat secara perlahan.}

\begin{itemize}
  \item \textbf{5. Peraturan Emas Sandaran: Prinsip Tiga-Dua-Satu}Ini adalah prinsip klasik dalam bidang keselamatan data, mudah diingat:
  \item \textbf{3}: Simpan sekurang-kurangnya tiga salinan data penting (satu asal, dua sandaran).
\end{itemize}

\booksection{2}

: Gunakan dua media penyimpanan yang berbeza untuk sandaran (contohnya satu di awan, satu di cakera keras luaran).

\begin{itemize}
  \item \textbf{1}: Sekurang-kurangnya satu sandaran diletakkan di lokasi fizikal yang berbeza (mengelakkan kebakaran, kecurian dan lain-lain yang memusnahkan semua data pada masa yang sama).
  \item \textbf{Untuk orang biasa, mencapai "satu di awan, satu di komputer" sudah dikira sangat baik. Tidak perlu terlalu rumit, yang penting anda benar-benar melakukannya.}Perhatian Khas untuk Rekod Sembang WeChat
  \item \textbf{Rekod sembang WeChat tidak boleh dipindahkan melalui sebarang alat pertukaran telefon pintar! Ia tidak diliputi oleh sandaran seluruh peranti. Mesti menggunakan fungsi "Pemindahan Rekod Sembang" terbina dalam WeChat, atau sandar ke WeChat versi komputer. Ramai orang hanya mendapati rekod sembang hilang sepenuhnya selepas bertukar ke telefon pintar baru, dan pada masa itu sudah terlambat — telefon pintar lama mungkin sudah dipulihkan kepada tetapan kilang.}Berlatih: Lakukan satu latihan sandaran lengkap
\end{itemize}

Periksa sama ada sandaran awan iCloud atau sandaran awan jenama telah diaktifkan, sahkan masa sandaran terakhir.

\begin{tcbwarn}{Sandar foto secara manual: Sambungkan komputer dengan kabel data, salin folder DCIM dan Pictures ke komputer.}
Sandar kenalan: Pengguna iPhone pergi ke iCloud.com eksport vCard; pengguna Android eksport fail .vcf dalam tetapan App Kenalan.

\end{tcbwarn}

\begin{tcbpractice}{Sandar rekod sembang WeChat: Buka WeChat Saya Tetapan Umum Pemindahan \& Sandaran Rekod Sembang Sandar ke komputer.}
\begin{enumerate}
  \item Salin satu salinan fail vcf kenalan yang dieksport dan folder foto ke cakera keras luaran atau pemacu USB, lengkapkan "satu salinan luar talian di tempat lain" dalam prinsip "Tiga-Dua-Satu".
  \item Bab 19: Pengesanan Kerosakan \& Penyelamatan Diri
  \item Telefon pintar seperti manusia, kadang-kadang juga "sakit". Lengah, keluar tiba-tiba, skrin hitam, tidak boleh dicas... kebanyakan masalah sebenarnya tidak setakut itu, anda boleh mencuba sendiri sebelum menghantar untuk pembaikan. Bab ini adalah "buku kecemasan telefon pintar" anda — cuba kaedah kos rendah terlebih dahulu, jika tidak berjaya baru cari doktor profesional.
  \item 1. Telefon Pintar Menjadi Lengah, Panas — Soalan yang Paling Sering Ditanya
  \item Jujur, telefon pintar mula menjadi lengah selepas digunakan kira-kira setahun, itu sangat biasa. Sebabnya tidak semestinya telefon pintar rosak, lebih kerap adalah "debu terkumpul" — fail sampah, proses latar belakang, ruang penyimpanan perlahan-lahan dipenuhi.
\end{enumerate}

\end{tcbpractice}
